%% Согласно ГОСТ Р 7.0.11-2011:
%% 5.3.3 В заключении диссертации излагают итоги выполненного исследования,
%% рекомендации, перспективы дальнейшей разработки темы.
%% 9.2.3 В заключении автореферата диссертации излагают итоги данного
%% исследования, рекомендации и перспективы дальнейшей разработки темы.

% =========================================================
% Рабочая структура заключения.
% Это НЕ финальная редакция, а удобный каркас для переноса и правки.
% Во всех блоках сохранены указания, откуда брать исходный материал:
%   - ВКР
%   - презентация
%   - рабочий план диссертации
% При последующей шлифовке служебные комментарии можно удалять.
% =========================================================

\section*{Основные научные результаты диссертации}
% Опора:
% ВКР: заключение, с. 71
% Презентация: слайды 25, 32, 38, 39
% План диссертации: блок "Заключение"
%
% Здесь удобно оформить компактным нумерованным списком:
% 1. результат по главе 1;
% 2. результат по главе 2;
% 3. результат по главе 3;
% 4. результат по главе 4.
%
% Желательно, чтобы формулировки были:
% - короткими;
% - проверяемыми;
% - максимально близкими к положениям, выносимым на защиту.
%
\begin{enumerate}
	\item %
	\item %
	\item %
	\item %
\end{enumerate}

\section*{Теоретическая значимость результатов}
% Опора:
% ВКР: введение, с. 9; заключение, с. 71
% Презентация: слайды 20--25, 27--32
% План диссертации: блок "Теоретическая значимость результатов"
%
% Здесь кратко суммировать:
% - исследование структуры модельного многообразия;
% - введение и интерпретацию точек оптимальной дальности;
% - построение аппроксимированного многообразия;
% - разработку линейных R-преобразований для перехода
%   к пространственной постановке.
%
% Черновой текст:

\section*{Практическая значимость результатов}
% Опора:
% ВКР: введение, с. 9; глава 5, с. 60--70; заключение, с. 71
% Презентация: слайды 34--38
% План диссертации: блок "Практическая значимость результатов"
%
% Здесь обязательно учесть:
% - получение начального приближения без решения полного семейства
%   краевых задач с нуля;
% - ускорение проектно-баллистического анализа;
% - применение для сценариев Земля--Марс и Земля--Венера;
% - программную реализацию в ILTT_Toolbox.
%
% Черновой текст:

\section*{Разработанные методы, алгоритмы и инструменты}
% Опора:
% ВКР: заключение, с. 71
% Презентация: слайды 10, 24, 32, 37--39
% План диссертации: блок "Разработанные методы, алгоритмы и инструменты"
%
% Здесь в концентрированном виде перечислить как единый подход:
% - методику построения, исследования и аппроксимации многообразия;
% - аппроксимированное многообразие плоской задачи;
% - линейные R-преобразования;
% - алгоритм выбора начального приближения;
% - программную реализацию ILTT_Toolbox;
% - при необходимости отдельно отметить зарегистрированную программу (РИД).
%
% Черновой текст:

\section*{Перспективы дальнейших исследований}
% Опора:
% ВКР: главы 3--5, с. 23--70
% Презентация: слайды 25, 32, 38
% План диссертации: блок "Перспективы дальнейших исследований"
%
% Здесь наметить:
% - расширение области аппроксимаций;
% - дальнейшее исследование малых витковых дальностей;
% - развитие пространственных преобразований;
% - перенос подхода на другие функционалы и постановки;
% - развитие программного инструментария для баллистического проектирования.
%
% Черновой текст:
